\begin{afterword}
\summaryhead

The post starts from the Ultimatum Game. Textbook decision theory says a responder should accept any positive offer, yet real responders refuse low ones, even when the stakes are two weeks' wages. The author sees an analogue in joining a group project: a would-be member holds out until the group changes something, such as how it invests its money, even though joining would still do more good than staying out.

It seems to the author that people in the ``atheist\slash{}libertarian\slash{}technophile/sf-fan/etcetera cluster'' set this price far too high. The explanation offered is instinct formed in bands of about 40 people, which underestimates the inertia of large groups, complex strategies and negotiations one does not see, together with weakness of will and a wish for status. The author admits speaking as an organizer.

The proposed rule: if joining still does net good and the issue is not worth fixing yourself, join anyway; pressure is justified only when a group blocks a competent member from fixing it. The post grants that leaders need some reason to answer to supporters, and ends by proposing this as a group norm.

\responsehead

The post is an exhortation built on an analogy and on the author's observations of a community, and it is candid about both. It says where the observations come from (the organizer's side, ``Perhaps I am a little prejudiced''), and it states the strongest objection to its own rule: that ``unconditional support would create its own problems.'' The rule also leaves room for dissent: holdouts are asked either to join or to do the fixing themselves, and pressure is allowed where a group blocks a competent member.

The evidence from the Ultimatum Game is stated a little more firmly than its source. The Wikipedia article the post links said that offers below 20\% are ``often'' rejected, and the post says ``most.'' A meta-analysis of 37 papers puts the average offer at 40\%. The point the post needs still holds: many people refuse low offers at a cost to themselves. The claim that ``this is just the way people feel'' needs one qualification. A study of fifteen small-scale societies found that play departs from the textbook in every society studied, but that how people play varies a great deal between societies.

The explanation has the least support. The post says it is ``almost certainly a matter of evolutionary psychology,'' and the only support given is that such situations would often have arisen among hunter-gatherers. The three mismatches that follow, large groups, complex strategies and unseen negotiations, are reasonable consequences of that account, but they have no more evidence than it does. The practical rule does not depend on the explanation, and would hold as advice without it.

The post asks whether its question has been studied. A related one had been: Albert Hirschman's \textsc{Exit, Voice, and Loyalty} (1970) analysed members who, disappointed by an organization, either leave or press for change from inside. Hirschman wrote about members rather than would-be joiners, but the post's advice, join and fix it from inside, resembles what Hirschman called voice.

In short: a candid exhortation to join groups more easily, whose rule of thumb names its own main objection and does not depend on the evolutionary explanation offered for it.
\end{afterword}
